% SPDX-License-Identifier: Apache-2.0
%
% PCIe in TxHDL. Built by //docs:pcie. Every listing is included from
% the tree, and every printout and figure is what the build made by
% running the example.
\documentclass[journal,9pt]{IEEEtran}

\newcommand{\listingsize}{\fontsize{6}{7}\selectfont}
% SPDX-License-Identifier: Apache-2.0
%
% The house preamble, shared by every document in this directory. Loaded
% after \documentclass, before \title. Keeping it in one file is what
% stops two articles drifting apart in font, listing style or figure
% style.
% Computer Modern. IEEEtran selects Times (ptm) by default, so the face has
% to be chosen explicitly.
%
% Latin Modern is the Computer Modern variety used here, and the choice is
% forced rather than aesthetic. Under T1 encoding, plain `cmr` has no Type 1
% outlines unless cm-super is installed, and it is not in the pinned
% distribution, so pdflatex falls back to EC bitmaps and every glyph in the
% PDF becomes a Type 3 font. That was measured: the first build produced a
% document with 10 Type 3 fonts and no Type 1. Latin Modern is Computer
% Modern redrawn with a full T1 repertoire and real .pfb outlines, so the
% same design comes out scalable.
\usepackage[T1]{fontenc}
\usepackage[utf8]{inputenc}
\usepackage{lmodern}
% microtype is not in the pinned TeX distribution. emergencystretch alone
% keeps the two-column measure from producing overfull lines.
\setlength{\emergencystretch}{3em}

\usepackage{amsmath}
\usepackage{amssymb}
\usepackage{amsthm}
\usepackage{booktabs}
\usepackage{array}
% enumitem is not in the pinned TeX distribution; the list spacing below
% does what \setlist would have done.
\usepackage{float}
\usepackage{xcolor}
\usepackage{listings}
\usepackage{tikz}
\usetikzlibrary{arrows.meta,positioning,fit,backgrounds,calc}
\usepackage[hidelinks,bookmarks=true]{hyperref}

% Numbered, definition-style box for a rule the reader is meant to apply.
\theoremstyle{definition}
\newtheorem{rulebox}{Rule}
\newtheorem{example}{Example}

% Unnumbered asides.
\theoremstyle{remark}
\newtheorem*{tip}{Tip}
\newtheorem*{pitfall}{Pitfall}

% Tighter lists, in place of enumitem's \setlist.
\setlength{\topsep}{2pt}
\setlength{\itemsep}{1pt}
\setlength{\parsep}{0pt}

\newcommand{\code}[1]{\texttt{#1}}
\newcommand{\term}[1]{\emph{#1}}

% Syntax highlighting. Four roles, distinguishable in colour and still
% distinguishable in greyscale, because an IEEE article gets printed: the
% keyword is the only bold role, the comment is the only italic one, and the
% two remaining roles differ in lightness as well as in hue.
\definecolor{kwblue}{RGB}{20,60,150}
\definecolor{cmtgreen}{RGB}{90,120,90}
\definecolor{strred}{RGB}{150,50,40}
\definecolor{numgray}{gray}{0.45}
\definecolor{framegray}{gray}{0.70}
\definecolor{bgshade}{gray}{0.97}
% A darker fill for figure cells. The listing background has to stay very
% light to keep code readable; a figure cell has to be clearly shaded or the
% distinction it draws is invisible in print.
\definecolor{cellshade}{gray}{0.90}

% The listing size is the document's choice, because it depends on the
% column: a two-column article sets `\listingsize` to 6pt and keeps its
% listings within 70 columns, a one-column document keeps the body size
% and includes source files at 80. Lines are never broken; a line that does not fit
% overflows the frame, which is visible, rather than wrapping, which is
% not.
\providecommand{\listingsize}{\scriptsize}

% `'` is deliberately not a string delimiter: the language uses it for loop
% labels, as Rust does, so treating it as a quote would swallow the rest of
% the line.
\lstdefinestyle{txhdl}{
  basicstyle=\ttfamily\listingsize,
  keywordstyle=\bfseries\color{kwblue},
  commentstyle=\itshape\color{cmtgreen},
  stringstyle=\color{strred},
  numberstyle=\tiny\color{numgray},
  morekeywords={mod,use,pub,const,type,struct,enum,trait,impl,fn,async,let,mut,
    match,if,else,loop,while,for,in,break,continue,return,as,await,
    module,bus,channel,transaction,protocol,pipeline,flow,seq,config,
    port,stage,pipe,instance,connect,bind,tag,domain,blackbox,foreign,
    property,role,signal,handler,true,false,
    sequence,interface,modport,implements,package,import,var,wire,func,proc,
    case,default,next,fork,branch,join,state,fsm},
  morecomment=[l]{//},
  morestring=[b]",
  showstringspaces=false,
  columns=fullflexible,
  keepspaces=true,
  breaklines=false,
  % Framed, so a listing is bounded rather than floating in the column.
  frame=single,
  rulecolor=\color{framegray},
  framesep=3pt,
  backgroundcolor=\color{bgshade},
  xleftmargin=3pt,
  xrightmargin=3pt,
  aboveskip=6pt,
  belowskip=6pt,
  % Every listing is numbered and captioned, so it can be referred to by
  % number. `caption` without `float` numbers the listing and leaves it
  % where it was written, which is what a listing in running prose needs.
  captionpos=b,
  abovecaptionskip=2pt,
  belowcaptionskip=6pt,
}

% Figures drawn as text. Framed the same way, and with the highlighting
% switched off, because a box-and-arrow diagram is not code and half its
% words turning blue reads as an accident. Setting a style adds keys rather
% than clearing the ones already set, so the three roles are neutralised by
% hand rather than by leaving them out.
\lstdefinestyle{ascii}{
  basicstyle=\ttfamily\listingsize,
  keywordstyle=\ttfamily,
  commentstyle=\ttfamily,
  stringstyle=\ttfamily,
  columns=fullflexible,
  keepspaces=true,
  frame=single,
  rulecolor=\color{framegray},
  framesep=3pt,
  backgroundcolor=\color{bgshade},
  xleftmargin=3pt,
  xrightmargin=3pt,
  aboveskip=6pt,
  belowskip=6pt,
}
\lstset{style=txhdl}

% House TikZ styles. `shorten` lives on the arrow styles rather than on each
% arrow, so every arrow in the document gets the gap, including ones added
% later. TikZ clips a path at a node's border, so without it an arrowhead
% butts against the box with no gap at all. Keep `node distance` well above
% twice the shorten, or the arrow shrinks to nothing.
\tikzset{
  box/.style={draw,rounded corners=2pt,align=center,inner sep=4pt,
              font=\footnotesize},
  boxf/.style={box,fill=bgshade},
  lbl/.style={font=\scriptsize\itshape,align=center},
  fl/.style={-{Stealth[length=4pt]},shorten >=2.5pt,shorten <=2.5pt},
  fld/.style={fl,dashed},
}



\InputIfFileExists{buildstamp.tex}{}{}
\providecommand{\buildstamp}{Draft build (outside the Bazel flow).}

% A range of a source file, by its begin/end markers.
\newcommand{\pinspart}[3]{%
\lstinputlisting[rangebeginprefix=//\ begin\{,rangebeginsuffix=\},%
rangeendprefix=//\ end\{,rangeendsuffix=\},includerangemarker=false,%
linerange=#1-#1,caption={#2},label={#3}]{lib/parts/src/bus/axi_pins.rs}}
\newcommand{\barpart}[3]{%
\lstinputlisting[rangebeginprefix=//\ begin\{,rangebeginsuffix=\},%
rangeendprefix=//\ end\{,rangeendsuffix=\},includerangemarker=false,%
linerange=#1-#1,caption={#2},label={#3}]{pcie/src/bar.rs}}

\title{PCIe in TxHDL}

\author{Filip Filmar and Dragiša Janković%
\thanks{Both authors are independent. Filip Filmar is the
corresponding author, at f@hdlfactory.com; Dragiša Janković is
reached at linkedin.com/in/dragisa-jankovic.}%
\thanks{This is the tutorial on putting TxHDL behind a PCIe endpoint,
for a reader who knows the AXI link~\cite{axi}. Every listing is
included from the project repository \code{HDL/txhdl} at build time,
and every printout and figure is what the build produced by running
the example. The cover~\cite{cover} lists every document.}%
\thanks{The authors used a large language model, Claude, as an
assistant in exploring the concepts and in writing and constructing
the documents and the programs; every commit of the repository says
so and carries the prompts that led to it, verbatim.}%
\thanks{\buildstamp}}

\markboth{PCIe in TxHDL}{Filmar and Janković: PCIe in TxHDL}

\begin{document}
\maketitle

\begin{abstract}
The Alinx AX7A200B has a PCIe x2 edge connector wired to the FPGA's
GTP transceivers. A PCIe endpoint sits on those transceivers and on
the FPGA's integrated PCIe block, and AMD provides it as an IP core.
The build generates it, the DMA/Bridge Subsystem for PCI Express at
Gen2 over both lanes with its bypass window open, and puts TxHDL
behind that window. A read or a write of the endpoint's BAR1 by a
host comes out of the core as a burst on an AXI4 master, at its pins. A new part, \code{AxiPins},
joins those pins to the AXI link's channels, and a lowered design of
four registers answers them. The part is simulated against its run
under \code{nvc} and Verilator, the design is driven through the pins
in Rust, and the whole goes through Vivado to a bitstream that meets
timing. It has not run in a PC, and what that needs is filed.
\end{abstract}

\section{What Is Here}
\label{sec:p-what}

Three things, in three places in the tree.

\code{AxiPins}, in \code{txhdl\_parts} under \code{bus::axi\_pins}, is
a host's AXI4 pins joined to the link's channels. It goes wherever a
host tracker would.

\code{pcie::bar}, under \code{//pcie}, is what BAR1 reaches:
\code{BarRegs}, four 64-bit words behind a peripheral tracker, and
\code{PcieBar}, one lowered unit that holds \code{AxiPins}, the
tracker and the registers.

\code{//pcie:xdma\_x2} is the endpoint, generated by Vivado, and
\code{pcie/board} is the hand-written top and the board's constraints.

\section{The Endpoint}
\label{sec:p-endpoint}

The build runs Vivado hermetically, and \code{rules\_vivado} has a rule
that generates an IP core from its name and its settings,
\code{vivado\_ip}. Its output is the core's sources and its
\code{.xci}; the synthesis rule reads that with \code{read\_ip} and
synthesises the core with the design.

Two of AMD's cores give a BAR as an AXI master on an Artix-7: the AXI
Memory Mapped to PCI Express bridge, \code{axi\_pcie}, and the
DMA/Bridge Subsystem, \code{xdma}. Both were tried at Gen2 over two
lanes, with the board's pins. Both placed on the
two GTP channels and the one integrated PCIe block with no error.
\code{axi\_pcie} missed timing by 0.097~ns on three paths, all inside
its AXI slave bridge, which nothing here uses and which it cannot leave
out. \code{xdma} can leave its slave out, and met timing. Its user
clock is 125~MHz: at Gen2 over two lanes the core refuses 62.5.

The core's settings are in \code{pcie/BUILD.bazel}: Gen2, two lanes,
the slot's 100~MHz reference clock, 64-bit data on the master, and
the bypass window open, a megabyte of memory at BAR1. On this part
the core has no bridge mode: a request for one is silently ignored,
and the core is a DMA engine, with its own registers at BAR0 and a
descriptor-driven master that nothing here drives. The bypass window
is the bridge: what a host reads or writes in BAR1 comes straight out
of the \code{m\_axib} master, with no descriptor in between. The
build asked for the bridge mode and got the DMA engine,
unnoticed until the root port simulation read BAR0~\cite{i451}. The
master's addresses are 64 bits wide; the design behind it takes the
low 32.

\section{The Pins}
\label{sec:p-pins}

A core from elsewhere that is an AXI4 host has pins: a \code{valid}
and a \code{ready} per channel, and a wire per field, named as the
specification names them. The parts of this library speak the link as
channels, one transaction a beat~\cite{axi}. \code{AxiPins}
stands between the two, and holds nothing.

\pinspart{ports}{The pins, as the unit's two sides.}{lst:p-ports}

The host's 24 pins are a struct of their own, \code{AxiHostPins},
which derives \code{Ports}, and the inputs hold them as one field,
\code{pins}, beside the two channels from the link; the unit's ports
are \code{pins\_awid} and the rest. A design that has a host's pins
among its own ports keeps them as one field too and passes it whole:
Vreteno's board holds its JTAG master's as \code{jtag}, whose ports
are \code{jtag\_awid} and the rest, and this design builds the field
from its own ports where it passes it (issue~\#579).

An address phase or a write beat the host offers goes onto its channel
in the step the channel has room, and the room is the host's
\code{ready}. A write response or a read beat at the head of its
channel is the host's \code{valid} and fields, and it is taken in the
step the host's \code{ready} is high. The burst type and the responses
are turned between their AXI4 encodings and the link's enums with
\code{select!}. The host drives no \code{AxQOS} and no
\code{AxREGION}, so both go out as zero.

\pinspart{unit}{\code{AxiPins}: the step.}{lst:p-unit}

The pins are wires, so a process that drives them and the unit that
reads them meet in one step. In a run, the host is therefore joined
ahead of the unit, so that the unit reads what the host drove in the
same step, as the netlist does.

\code{ex\_axi\_pins} puts a host on the pins, written as a simulation,
in front of \code{AxiPins<32, 64, 8, 4>} with a peripheral tracker and
a RAM behind: the endpoint's widths. It writes three 64-bit words, reads
four back, and reads past the RAM's end, which the peripheral refuses.

\lstinputlisting[caption={\code{ex\_axi\_pins.rs}: a host on pins, and a
RAM behind them.},label={lst:p-ex}]{lib/examples/ex_axi_pins.rs}

\lstinputlisting[style=ascii,caption={What \code{ex\_axi\_pins} printed
when the build ran it.},label={lst:p-out}]{out_axi_pins.txt}

\begin{figure}[htbp]
\centering
\input{axi_pins_timing.tex}
\caption{The write through the pins: its address phase and its three
beats handed to the channels as the host offers them, and its response
handed back.}
\label{fig:p-pins}
\end{figure}

The unit is lowered, and its netlist agrees with this run under
\code{nvc} and under Verilator: \code{//docs:axi\_pins\_sim\_axi\_pins\_tb\_test}
and \code{//docs:axi\_pins\_vsim\_test}.

\section{Behind BAR1}
\label{sec:p-bar}

\code{BarRegs} is four words, at the offsets a host reads and writes in
BAR1.

\begin{center}\footnotesize
\begin{tabular}{@{}llp{0.55\columnwidth}@{}}
\toprule
Offset & Word & Access \\
\midrule
\code{0x00} & Identifier & Read only: \code{TxHDL} in ASCII and a
version, \code{0x5478\_4844\_4c00\_0001}. \\
\code{0x08} & Scratch & Read and write. \\
\code{0x10} & LEDs & Read and write; the low two bits light LED3 and
LED4. \\
\code{0x18} & Writes & Read only: writes served, of any word. \\
\bottomrule
\end{tabular}
\end{center}

It takes one request at a time, as Vreteno's timer does: a read is
answered in the cycle it is taken, and a write when its beat comes. It
looks at bits 4 and 3 of an address alone, since the tracker sends it
only its own range, and a write replaces the whole word whatever its
strobe.

\barpart{regs}{\code{BarRegs}: four words behind a peripheral
tracker.}{lst:p-regs}

\code{PcieBar} is the three units as one: the pins, the tracker, and the
registers, joined by channels. Its ports are the endpoint's master
pins, a reset and the LEDs, so the board's top wires the core to it
pin for pin.

\barpart{bar}{\code{PcieBar}: everything behind BAR1, as one
unit.}{lst:p-bar}

\code{//pcie:bar\_test} drives \code{PcieBar} through the pins as the
endpoint does when a host reads and writes the BAR. It reads the
identifier, writes the scratch word and reads it back with its
identifier, lights an LED, writes the identifier and finds it
unchanged, and reads the count of writes. It also checks that the
netlist is one module with the pins as its ports.

\section{The Board}
\label{sec:p-board}

The top is written by hand, and holds what a lowered unit cannot say:
the buffer for the transceivers' reference clock, the core, and the
core's clock and reset handed to the lowered design. \code{PcieBar}
runs on the core's 125~MHz user clock, so nothing crosses a clock, and
its reset is high while the core's is low.

\lstinputlisting[caption={\code{pcie/board/pcie\_top.v}: the
endpoint and the design behind it.},label={lst:p-top}]{pcie/board/pcie_top.v}

The constraints name the two lanes and the reference clock from the
table in the board's manual, the reset key, and the LEDs. A
transceiver's pins take no I/O standard. The slot's reset, PERST\#, is
not in the manual's table~\cite{i179}, so the core is reset from the
board's reset key, as the other board designs are.

\lstinputlisting[caption={\code{pcie/board/ax7a200\_pcie.xdc}: the
board's pins.},label={lst:p-xdc}]{pcie/board/ax7a200_pcie.xdc}

\code{bazel build //pcie:endpoint\_synth} generates the core, lowers
\code{PcieBar}, and puts both and the top through Vivado's synthesis
for the board's part, with no error and no critical warning; the
warnings are about logic inside the core that its settings leave
unused. \code{bazel build //pcie:endpoint\_pnr} places and routes the
design with the board's pins and writes the bitstream, in about eight
minutes. Every timing constraint is met, with 0.656~ns of setup slack
and 0.036~ns of hold slack over 42\,813 endpoints, on the reference
clock and the clocks the core derives from it. No register is without
a clock and no internal endpoint is unconstrained. The design rule
check finds two warnings, both inside the core: a LUT input its
equation does not use, and 23 nets with no load. The design takes
11\,320 LUTs, 13\,570 flip-flops, 20.5 block RAM tiles, the two GTP
channels, the integrated PCIe block, and five pins besides the
transceivers'. The core accounts for nearly all of it.

\section{What Is Not Done}
\label{sec:p-not}

Nothing here has run in a PC. That is the test that says the link
trains at Gen2 over two lanes, that the host enumerates the device and
its BAR, and that a read and a write reach the registers; it needs the
board in a slot and a small program on the host~\cite{i180}.

The endpoint is simulated, on its own and against a root port. The
core's simulation library holds every source the core needs, since
the build compiles it from Vivado's own exported simulation rather
than from the top wrapper alone. \code{//pcie/sim:xdma\_test}
elaborates the board's top, the core inside it and \code{BarRegs}
behind BAR1 under Vivado's simulator and runs them with nothing on
the lanes: the link stays down, read as a known zero rather than an
unknown. \code{//pcie/sim:rp\_test} puts the same top on two lanes
against the behavioural root port and the host that Vivado writes
into the core's example design, which the build keeps beside the
core~\cite{i178}. The host trains the link, at Gen2 over both lanes,
reads the device and vendor identifiers, scans and programs the BARs,
and then runs \code{pcie/sim/rp/tests.vh}: it reads the two words of
the identifier, writes the scratch word and reads it back, and reads
the count of writes, expecting what \code{BarRegs} keeps. It takes
twenty-six minutes and passes. Its first run, with the core as it
was, read BAR0 and found 64~KiB whose first word was the core's
own, \code{0x1fc0\_0006}: the request for the bridge mode had been
ignored and the core was a DMA engine~\cite{i451}. The bypass window
is the answer, above.

\code{vivado\_ip} reports success when Vivado refuses a core's
settings, so a mistyped setting shows only in the log~\cite{i177}. The
slot's PERST\# is not wired, as above~\cite{i179}.

\code{PcieBar} is its own bus, on the endpoint's clock. Putting the BAR
on Vreteno's bus instead is a crossing between two clocks, which
\code{ChanCdc} writes in the language and the lowering co-simulates
since issue 131~\cite{i131}; it is not done here, and the endpoint
keeps its own bus. \code{BarRegs} answers one
beat a burst and ignores a write's strobe, which a host reading and
writing whole registers does not notice.

\begin{thebibliography}{9}

\bibitem{axi}
F.~Filmar and D.~Janković, ``AXI in TxHDL,'' project repository
\code{HDL/txhdl}, \code{//docs:axi}, 2026.

\bibitem{i131}
Issue 131, ``feat(test): co-simulate a unit of several clocks, and
lower a clock-crossing channel,'' project repository \code{HDL/txhdl},
2026.

\bibitem{i177}
Issue 177, ``fix(rules\_vivado): vivado\_ip succeeds when Vivado
rejects the IP's configuration,'' project repository
\code{HDL/txhdl}, 2026.

\bibitem{i178}
Issue 178, ``feat(pcie): simulate the PCIe endpoint, its IP and a root
port,'' project repository \code{HDL/txhdl}, 2026.

\bibitem{i179}
Issue 179, ``board: the AX7A200B's PCIe reset, PERST\#, and its pin,''
project repository \code{HDL/txhdl}, 2026.

\bibitem{i180}
Issue 180, ``test(pcie): run the endpoint in a PC and read its BAR from
the host,'' project repository \code{HDL/txhdl}, 2026.

\bibitem{i451}
Issue 451, ``fix(pcie): the endpoint is a DMA engine, not a bridge;
the Artix-7 XDMA ignores functional\_mode,'' project repository
\code{HDL/txhdl}, 2026.

\bibitem{cover}
F.~Filmar and D.~Janković, ``TxHDL documents,'' project repository
\code{HDL/txhdl}, \code{//docs:cover}, 2026.

\end{thebibliography}

\end{document}
